%==============================================================================
% RESUMO (NBR 6028 — 150 a 500 palavras)
%==============================================================================

\begin{resumo}

Este livro oferece uma abordagem completa e autodidata sobre a interseção entre Odontologia e Inteligência Artificial. Organizado em 5 partes e 28 capítulos, totalizando aproximadamente 500 laudas, a obra conduz o leitor do nível zero (leigo) ao nível PhD/desenvolvedor. A Parte I apresenta os fundamentos históricos da IA na odontologia, introdução ao machine learning e Python com Google Colab. A Parte II aprofunda a teoria de redes neurais convolucionais, segmentação semântica, radiologia odontológica com deep learning e classificação de lesões bucais e câncer oral. A Parte III estabelece a metodologia de engenharia de software odontológico com Spec-Driven Development (SDD) e Test-Driven Development (TDD), incluindo projetos guiados, interpretabilidade de modelos, auditoria de código e reprodutibilidade científica. A Parte IV explora aplicações avançadas em periodontia, implantodontia, ortodontia, gêmeos digitais, modelos fundacionais e RAG odontológico. A Parte V aborda publicação científica, ética, LGPD, regulação, educação odontológica com IA e perspectivas futuras. Cada capítulo combina revisão bibliográfica rigorosa com citações auditáveis (DOI ativo), prática orientada por especificação (SDD) e testes (TDD) em Python e Google Colab. A obra inclui mais de 100 artigos científicos referenciados, 93 testes automatizados, fichamentos com resenha crítica e datasets odontológicos públicos.

\textbf{Palavras-chave:} Odontologia. Inteligência Artificial. Deep Learning. Processamento de Imagens. Gêmeos Digitais. SDD. TDD. Python. Google Colab.
\end{resumo}
